\subsection{The exact ordinary capitulation kernel and three remaining S-class tests}

\begin{statement}{E.84}{Proved}{Normalizer action evaluates the ordinary transfer and reduces the S-unit problem}
Use the actual fields and prime orderings of E.81:
\[
 \begin{gathered}
 K=\mathbb Q(\zeta),\quad \xi^3=\zeta,\quad a^3=17,\quad
 \beta_1=a\xi,\quad \beta_2^3=\beta_1,\\
 F=B_1=K(\xi,a),\qquad E=B_2=F(\beta_2),\qquad
 G=\operatorname{Gal}(E/F)=\langle\sigma\rangle .
 \end{gathered}
 \tag{E84.1}
\]
Write \(Q_1,Q_2,Q_3\) for the upper primes above three and
\(T_1,T_2,T_3\) for those above seventeen. Their corresponding
base classes are respectively \((1,0)\) and \((2,0)\) in the
certified group \(\operatorname{Cl}(F)=(\mathbb Z/3)^2\).
Let \(P=(73,t-33)\), \(t=a-\xi\), and \(I=P\mathcal O_E\);
then \([P]=(2,1)\).

The ordinary ideal-extension map is completely evaluated:
\[
 \begin{aligned}
 j:\operatorname{Cl}(F)&\longrightarrow\operatorname{Cl}(E)\{3\},
       &j(u,v)&=v[I],\\
 \ker j&=\langle(1,0)\rangle\simeq\mathbb Z/3,
       &\operatorname{im}j&=\langle[I]\rangle\simeq\mathbb Z/3.
 \end{aligned}
 \tag{E84.2}
\]
Every one of the six ramified prime classes has exact order
three. In addition,
\[
 [Q_1]+[Q_2]+[Q_3]=0,\qquad
 [T_1]+[T_2]+[T_3]=0.
 \tag{E84.3}
\]
Thus all the ideals \(Q_i^3,T_i^3,Q_1Q_2Q_3\) are principal;
the second product in (E84.3) is literally \((\beta_2)\).

Here is the exact remaining S-divisor map. Put
\[
 \begin{gathered}
 S=(Q_1,Q_2,Q_3,T_1,T_2,T_3),\quad
 \Delta=(-1,1,0,0,0,0)^{\mathsf t},\\
 L_0=3\mathbb Z^6+
       \mathbb Z(1,1,1,0,0,0)^{\mathsf t}
       +\mathbb Z(0,0,0,1,1,1)^{\mathsf t},\\
 \delta=[Q_2]-[Q_1],\qquad
 w=[I]+[T_2]-[T_1].
 \end{gathered}
 \tag{E84.4}
\]
The lattice \(L_0\) consists of principal S-divisors and has
index \(81\). Its HNF and quotient are
\[
 H_0=
 \begin{pmatrix}
 3&0&1&0&0&0\\0&3&1&0&0&0\\0&0&1&0&0&0\\
 0&0&0&3&0&1\\0&0&0&0&3&1\\0&0&0&0&0&1
 \end{pmatrix},
 \qquad \mathbb Z^6/L_0\simeq(\mathbb Z/3)^4.
 \tag{E84.5}
\]
Let \(\phi\) be the actual genus Artin map of E.81 and
\[
 W=\ker\!\left(\phi:
       \operatorname{Cl}(E)\{3\}^{G}\longrightarrow\Gamma\right).
 \tag{E84.6}
\]
Both \(w\) and \(\delta\) belong to \(W\), whose order divides
three. On \(\mathbb Z^6/L_0\), the kernel of the genus map is
exactly the line generated by \(\Delta\). Consequently the
full lattice \(L\) of principal S-divisors satisfies the
unconditional exact identity
\[
 L/L_0=
 \ker\!\left(\mathbb F_3\xrightarrow{\ t\mapsto t\delta\ }W\right)
 \quad\text{inside}\quad
 \mathbb Z^6/L_0,\qquad t\longmapsto t\Delta .
 \tag{E84.7}
\]
The image of the six S-prime classes still has order \(27\)
or \(81\); (E84.7) identifies its entire remaining kernel.

The twenty-seven ideals of E.81 now have the exact class formula
\[
 [J_{a,b,c}]=w+(b+2c)\delta .
 \tag{E84.8}
\]
Thus the S-class of \(I\) vanishes exactly when one of the three
literal ideals
\[
       I\,T_2T_1^{-1}(Q_2Q_1^{-1})^t,
                    \qquad t=-1,0,1,
 \tag{E84.9}
\]
is principal. Each of these three class values occurs nine
times in the old list. The values of \(w,\delta\) in \(W\)
are not evaluated here.

There is also an evaluated comparison with the original
complex tower. The fixed field of the automorphism \(h\)
below is \(L_2=\mathbb Q(\beta_2)\); its discriminant is
\(-3^{45}17^{16}\), and \(E/L_2\) is unramified everywhere
and split at the two primes above \(3,17\). The exact
ordinary transfer gives
\[
 \operatorname{ord}[\mathfrak q_{3,2}]=3,\qquad
 j_{E/L_2}[\mathfrak q_{3,2}]
       =[Q_1Q_2Q_3]=0.
 \tag{E84.10}
\]

Finally retain the exact distinction between the two norm
maps that arise in the proposed unit calculation. For
\(c_1=\beta_1+1+3\zeta\),
\[
 N_{F/K(\xi)}c_1=1-\zeta,\qquad
 c_1\notin N_{E/F}E^\times,\qquad
 \xi c_1\in N_{E/F}E^\times .
 \tag{E84.11}
\]
The last membership is the proved field-norm assertion of
E.79. Every S-unit preimage of \(\xi c_1\) has the now
proved principal ideal \(Q_1Q_2Q_3\). Choose a generator
\(\gamma\) of that ideal. The entire remaining norm fiber is
the explicitly translated ordinary-unit fiber
\[
 \{x\in\mathcal O_E[1/3,1/17]^\times:N_{E/F}x=\xi c_1\}
 =
 \gamma\,
 \{u\in\mathcal O_E^\times:
            N_{E/F}u=\xi c_1/N_{E/F}\gamma\}.
 \tag{E84.12}
\]
The quotient on the right is a unit of \(F\), and changing
\(\gamma\) changes it by an actual unit norm. Neither this
unit-norm coset nor the eventual exponent \(h\) of E.69 is
evaluated by the ordinary capitulation calculation.
The manuscript anchor is LM~\S51.16, fols.~733v--736v,
scans~172--178.
\end{statement}

\begin{proof}
The essential extra input is an actual normalizer action,
including its action on the certified base class group.
Define
\[
 \begin{array}{c|cc|cc}
       &\xi&\beta_2&a&\beta_1\\ \hline
 h&\zeta\xi&\beta_2&\zeta^2a&\beta_1\\
 g&\xi&\xi\beta_2&\zeta a&\zeta\beta_1
 \end{array}
 \qquad
 h^3=1,\quad g^3=\sigma,\quad g^9=1,\quad
 hgh^{-1}=g^4 .
 \tag{E84.13}
\]
These are automorphisms of the displayed field, not
formal actions on a class group. For example
\(\beta_2^9=17\xi^3\); both substitutions preserve this
equation and \(\xi^6+\xi^3+1=0\), and their inverse
substitutions are powers of themselves. Their \(27\)
distinct products \(g^i h^j\), \(0\le i<9,0\le j<3\),
are all the \(K\)-automorphisms, since \([E:K]=27\).
Both normalize \(F\) and commute with \(\sigma=g^3\).

In the actual six-prime ordering the permutations are
\[
 h=(Q_1\,Q_2\,Q_3)(T_1\,T_2\,T_3),
 \qquad
 g=(Q_1\,Q_3\,Q_2),\qquad g(T_i)=T_i.
 \tag{E84.14}
\]
The certificate checks these as equalities of ideals in
the maximal order of \(E\), using the substitutions
(E84.13). In the primitive base coordinate \(t=a-\xi\),
the two substitutions are
\(t\mapsto\zeta^2a-\zeta\xi\) and
\(t\mapsto\zeta a-\xi\). Their actions on the certified
class basis of E.76 are
\[
 h_F=
 \begin{pmatrix}1&2\\0&1\end{pmatrix},\qquad
 g_F=
 \begin{pmatrix}1&1\\0&1\end{pmatrix}
 \quad\text{over }\mathbb F_3.
 \tag{E84.15}
\]
For each basis ideal and for \(P\), the program retains
the actual image ideal, its class coordinates and a
literal principal generator of its quotient by the
specified class-basis ideals. A separate rational
matrix calculation verifies these six principal-ideal
identities in the complete integral basis of \(F\).
Thus (E84.15) uses the exact certified class map, with
the arithmetic automorphisms fixed by (E84.13)
\cite[entries \texttt{bnfisprincipal} and
\texttt{nfgaloisapply}]{E_PARICertifiedUnits}.

We spell out why these automorphisms preserve the
particular kernel \(W\). Let \(\mathcal H_E\) be the
maximal unramified abelian three-extension of \(E\),
and let \(\mathcal G\) be its maximal subfield which
is abelian over \(F\), as constructed in E.76.
An extension to an algebraic closure of an automorphism
normalizing \(E\) and \(F\) carries
\(\mathcal H_E\) to itself: unramifiedness and
abelianity are preserved by conjugation, and the
inverse gives the reverse inclusion. It carries
\(\mathcal G\) to itself for the same reason and its
maximal property. The Artin map is equivariant:
on an unramified prime this is the equality
\[
 \operatorname{Frob}_{\nu\mathfrak P}
       =\nu\operatorname{Frob}_{\mathfrak P}\nu^{-1},
       \qquad \nu=g,h,
 \tag{E84.16}
\]
and multiplication gives the assertion on every ideal
class. Since \(g,h\) commute with \(\sigma\), they
also preserve the subgroup of \(\sigma\)-fixed
classes. Hence both preserve the kernel in (E84.6).
The retained ray matrices additionally verify that
they preserve the exact relation lattice of E.76,
and that their genus actions carry each of its six
Artin columns to the column specified by (E84.14).

E.81 proves \(|W|\mid3\). An automorphism of a group
of order three has order dividing two. The actions
of \(g,h\) have orders dividing nine and three,
respectively, so both are the identity on \(W\).
This argument also includes \(W=0\).

Put \(k=j(1,0)\), \(i=[I]\), and \(t_j=[T_j]\).
The actual Artin matrix of E.81 gives
\(\phi(i+t_2-t_1)=0\), so \(w=i+t_2-t_1\in W\).
By (E84.15),
\[
 g(i)-i=j\bigl((0,1)-(2,1)\bigr)=j(1,0)=k.
 \tag{E84.17}
\]
Equation (E84.14) fixes both \(t_1,t_2\). Since \(g\)
acts trivially on \(W\), (E84.17) implies
\(k=(g-1)w=0\).
On the other hand E.81 proves that \(i\) has order
exactly three, through its nonzero Artin value
\((6,0,0,0)\) in \(\mathbb Z/9\oplus(\mathbb Z/3)^3\).
Because \((2,1)=2(1,0)+(0,1)\), this proves
(E84.2) on the entire certified base class group.

Extension of the six base primes gives exactly
their upper cubes. Thus
\[
 3[Q_j]=j(1,0)=0,\qquad
 3[T_j]=j(2,0)=0.
 \tag{E84.18}
\]
Every one of their Artin columns in E.81 is nonzero;
their orders are therefore exactly three.
The three \(Q\)-columns are equal, so
\(\delta=[Q_2]-[Q_1]\) belongs to \(W\).
The identity \(h\delta=\delta\), together with the
actual cycle (E84.14), gives
\([Q_3]-[Q_2]=\delta\). Therefore
\([Q_3]=[Q_1]+2\delta\), and their sum is
\(3[Q_1]+3\delta=0\). The \(T\)-product is
\((\beta_2)\), since this element is a uniformizer
at each of the three primes over seventeen and a
unit at every other finite prime. This proves
(E84.3) and all the principal relations defining
\(L_0\).

The displayed HNF has determinant \(81\); its full
integral Smith maps and inverse maps are retained.
The genus images of the six S-primes have order
\(27\) by E.81. Since \(\Delta\notin L_0\) has
order three and has zero Artin image, it is the
entire genus kernel on \(\mathbb Z^6/L_0\).
Every principal divisor has zero Artin image.
On that remaining line the actual class map is
exactly \(t\Delta\mapsto t\delta\), proving (E84.7)
in both directions.

For the comparison of the old list with (E84.9),
substitute
\([Q_2]=[Q_1]+\delta\) and
\([Q_3]=[Q_1]+2\delta\) into its literal exponent
vector \((3a+2b+2c,b,c,-1,1,0)\).
This gives (E84.8). There is an exact ideal
comparison behind the reduction. Choose the centered
representative \(t\) of \(b+2c\) modulo three and
write \(m=(b+2c-t)/3\). The difference of its
exponent vector from the vector of (E84.9) is
\[
 c(1,1,1,0,0,0)
 +3(a+b+c-m)(1,0,0,0,0,0)
 +3(m-c)(0,1,0,0,0,0).
 \tag{E84.19}
\]
All three displayed factors have just been proved
principal. Thus multiplying by their principal
functions gives the actual ideal-class comparison;
no identification of abstract vector spaces is
substituted for it. The linear map
\((a,b,c)\mapsto b+2c\) has kernel of order nine,
which proves the multiplicity assertion.

The fixed field of \(h\) contains \(\beta_2\) and
has degree eighteen. This is exactly
\(L_2=\mathbb Q(\beta_2)\), with polynomial
\(X^{18}+17X^9+289\). The maximal-order calculation
gives its discriminant \(-3^{45}17^{16}\);
its cube equals the independently retained
discriminant \(-3^{135}17^{48}\) of \(E\).
The relative discriminant ideal therefore has
norm one. All infinite places are complex, so the
extension is unramified everywhere. The two
three-cycles in (E84.14) give complete splitting
at \(3,17\), also agreeing with the all-level
comparison in E.76. The original prime extends
to \(Q_1Q_2Q_3\), which is principal by (E84.3).
E.72 proves that its class in \(L_2\) has exact
order three. This establishes the concrete
capitulating class (E84.10).

For the norm comparison, put \(u=1+3\zeta\).
The literal identities
\(u^3=1-18\zeta\) and \(\beta_1^3=17\zeta\)
give \(N_{F/K(\xi)}(\beta_1+u)=1-\zeta\).
At each base prime over seventeen the residue
field has order \(289\), and the cubic relative
extension is totally ramified. The residue of a
unit norm is therefore a cube. The certificate
checks, in each of the three actual residue maps,
\[
  c_1^{96}=\zeta,\qquad
  (\xi c_1)^{96}=1 .
 \tag{E84.20}
\]
Since \(\zeta\ne1\), the first identity excludes
the uncorrected norm. The second is the tame
part of the corrected field-norm calculation;
E.79 supplies its full six-place verification and
the Hasse norm theorem, yielding (E84.11).

The actual base valuations of \(\xi c_1\) are
\((1,1,1,0,0,0)\). Each of the six relative
extensions has residue degree one and a single
upper prime, so ideal norm preserves these
valuation coefficients. An S-unit preimage
therefore has divisor precisely
\(Q_1+Q_2+Q_3\). All elements with that divisor
are \(\gamma\mathcal O_E^\times\), and taking
their literal norms gives (E84.12). The ratio
there is a unit because its divisor is zero.
Changing \(\gamma\) by an ordinary unit multiplies
the ratio by the inverse of its norm, proving
the stated independence of the remaining coset.

The searches retained with this result found no
small generator or norm preimage in their stated
finite search sets. They are not used to infer
nonprincipality or failure of an S-unit norm.
In particular the principality conclusions
above follow from the exact normalizer and
class-group argument, even where those searches
did not produce a principal function.
\end{proof}
